% ==============================================================================
% TEMPLATE ACUAN MATA KULIAH: PPG (PEMROGRAMAN PERMAINAN / GAME)
% ==============================================================================
% Dokumen ini disusun persis sesuai format berkas referensi asli PPG (tanpa preamble).
% Bab utama langsung diawali dari \chapter{Dokumentasi Praktikum}.
%
% ATURAN MUTLAK FORMAT PPG:
% 1. Laporan PPG TIDAK MEMILIKI bab Dasar Teori dan TIDAK MEMILIKI bab Kesimpulan.
% 2. Bab utama adalah \chapter{Dokumentasi Praktikum} (bukan Hasil dan Pembahasan).
% 3. Wajib memisahkan \section{Hasil Praktikum} dan \section{Hasil Akhir} dengan \clearpage.
% 4. Tata letak gambar SATU GAYA TUNGGAL:
%    - \begin{figure}[H] lebar selalu 0.6\textwidth (atau 0.4\textwidth untuk panel sidebar)
%    - Tepat di bawahnya: 'Keterangan Gambar :'
%    - Diikuti \begin{itemize} (Pola 1 deskriptif atau Pola 2 Cara+Fungsi).
% ==============================================================================

\chapter{Dokumentasi Praktikum}
\label{chap:dokumentasipraktikum}

% ==============================================================================
% 1. SEKSI HASIL PRAKTIKUM (PROSES DAN LANGKAH KERJA)
% ==============================================================================
\section{Hasil Praktikum}
\label{sec:hasilpraktikum}

\subsection{Inisialisasi Lingkungan dan Pembuatan Proyek Game Engine}

\begin{figure}[H]
    \centering
    \includegraphics[width=0.6\textwidth]{1.png}
    \caption{Pembuatan Proyek Baru Menggunakan Template pada Game Engine}
    \label{fig:1}
\end{figure}
Keterangan Gambar :
\begin{itemize}
    \item Buka aplikasi editor game engine dan pilih kategori template proyek yang sesuai (\textit{Third Person} atau \textit{First Person}).
    \item Tentukan direktori penyimpanan lokal, beri nama proyek sesuai format penamaan modul, lalu klik tombol pembuatan proyek.
\end{itemize}

\subsection{Eksplorasi Antarmuka Editor dan Hierarki Objek}

\begin{figure}[H]
    \centering
    \includegraphics[width=0.4\textwidth]{2.png}
    \caption{Panel Navigasi Objek dan Properti pada Editor Game Engine}
    \label{fig:2}
\end{figure}
Penjelasan Antarmuka :
\begin{itemize}
    \item Gunakan panel \textbf{Viewport} utama untuk menavigasikan kamera dan menata objek 3D di dalam level lingkungan permainan.
    \item Tinjau panel hierarki (\textbf{Outliner}) untuk memilih aktor dan panel \textbf{Details} untuk mengubah atribut parameter fisik.
\end{itemize}

\subsection{Implementasi Logika Visual Scripting (Blueprint)}

% Contoh Pola 2: Terstruktur Cara + Fungsi
\begin{figure}[H]
    \centering
    \includegraphics[width=0.6\textwidth]{3.png}
    \caption{Penyusunan Rangkaian Node Logika Interaksi Karakter}
    \label{fig:3}
\end{figure}
Penjelasan :
\begin{itemize}
    \item \textbf{Cara:} Klik kanan pada kanvas graf logika visual (\textit{Event Graph}), panggil fungsi eksekusi yang ditargetkan, lalu hubungkan pin eksekusi (\textit{white execution line}) ke variabel tujuan.
    \item \textbf{Fungsi:} Berfungsi untuk memvalidasi pemicu interaksi (\textit{trigger collision}) saat karakter menyentuh area objek dalam permainan.
\end{itemize}

% ==============================================================================
% 2. SEKSI HASIL AKHIR (VERIFIKASI FUNGSIONAL GAMEPLAY)
% Wajib diawali \clearpage
% ==============================================================================
\clearpage
\section{Hasil Akhir}
\label{sec:hasilakhir}

Bagian ini memaparkan hasil verifikasi fungsional gameplay, animasi, maupun modifikasi perbaikan fitur yang telah diimplementasikan pada sesi praktikum mandiri.

\subsection{Pengujian Mekanika Fisika dan Navigasi Level}

\begin{figure}[H]
    \centering
    \includegraphics[width=0.6\textwidth]{4.png}
    \caption{Uji Coba Fungsionalitas Navigasi Karakter di Area Tanjakan dan Rintangan}
    \label{fig:4}
\end{figure}
Keterangan Gambar :
\begin{itemize}
    \item Karakter pemain dapat bergerak secara stabil melintasi rintangan tanpa terjadi galat tabrakan objek (\textit{collision clipping}).
    \item Sistem kamera pemain mengikuti perpindahan arah karakter secara halus dan konsisten terhadap orientasi kontrol gerakan.
\end{itemize}

\subsection{Verifikasi Siklus Kondisi Karakter dan Interaksi Lingkungan}

\begin{figure}[H]
    \centering
    \includegraphics[width=0.6\textwidth]{5.png}
    \caption{Kondisi Respon Status Karakter Saat Memasuki Zona Target}
    \label{fig:5}
\end{figure}
Penjelasan :
\begin{itemize}
    \item \textbf{Deskripsi:} Pengujian respon sistem saat karakter menyentuh zona akhir penugasan.
    \item \textbf{Keterangan:} Indikator grafis dan efek visual terpicu secara simultan, membuktikan bahwa logika algoritma pemicu status telah bekerja 100\% fungsional tanpa menimbulkan \textit{bug} eksekusi.
\end{itemize}

% ==============================================================================
% 3. DAFTAR PUSTAKA
% ==============================================================================
\chapter{Daftar Pustaka}
\label{chap:daftarpustaka}

\begin{enumerate}
    \item \label{modulPPG} Tim Asisten dan Dosen Pengampu. \textit{Modul Praktikum Pemrograman Permainan: [Topik Pertemuan X]}. Yogyakarta: Departemen Teknik Elektro dan Informatika, Sekolah Vokasi, Universitas Gadjah Mada, 2026.
    
    \item \label{gregory} J. Gregory. \textit{Game Engine Architecture}, 3rd ed. Boca Raton, FL: CRC Press, 2018.
    
    \item \label{romero} M. Romero. \textit{Blueprints Visual Scripting for Unreal Engine 5}, 3rd ed. Birmingham, UK: Packt Publishing, 2022.
    
    \item \label{akenine} T. Akenine-Möller, E. Haines, dan N. Hoffman. \textit{Real-Time Rendering}, 4th ed. Boca Raton, FL: A K Peters/CRC Press, 2018.
    
    \item \label{engineDoc} Epic Games. \textit{Unreal Engine 5 Documentation}. [Online]. Tersedia: \url{https://docs.unrealengine.com/}.
    
    \item \label{sourcecode} [Nama Mahasiswa], \textit{S5-PPG-[NIM] (Source Code Proyek Game Praktikum)}. [Online]. Tersedia: \url{https://github.com/[Username]/[Repo-Game]}.
\end{enumerate}
